\chapter{Engineering Design and Analysis}
\label{ch:engineering_design}

\section{Airfoil and wing analogous estimation}
According to \textbf{HUG-OPS-005}, HUGO is required to operate at speeds up to Mach 0.75\cite{BaselineReport}. This is the flight condition at which tests are expected to be conducted during operations. In this section, an estimation of the lift characteristics of HUGO is performed based on statistics from suitable reference aircraft, along with a preliminary airfoil selection suitable for the considered operations.


\subsection{Reference aircraft parameters}
\label{sub:lift_reference_aircraft}
The initial step in selecting an airfoil is to determine the design lift coefficient $C_{L_\text{des}}$ and maximum lift coefficient $C_{L_\text{max}}$. Several steps have to be conducted beforehand to obtain these coefficients. Due to the specified testing conditions, $C_{L_\text{des}}$ was selected for the Mach 0.75 cruise airspeed condition at FL270 (8230 m). For initial estimates, the NASA DAST project \footnote{\label{fn:DAST}\href{https://web.archive.org/web/20230310084826/https://www.nasa.gov/centers/dryden/multimedia/imagegallery/DAST/DAST_proj_desc.html}{NASA - Firebee Drone: DAST} (accessed 07.05.2026)}, NASA X-56A UAS \footnote{\href{https://www.nasa.gov/aeronautics/x-56a-mutt/}{NASA X-56 UAS testbed} (accesed on 8.5.2026)} and the FLEXOP UAS \cite{Yu2022} are used as reference aircraft.\\
DAST was initially developed as a modified version of the Teledyne-Ryan BQM-34E/F Firebee II, with the purpose of remotely testing new aerodynamic designs in the 1980s. The aircraft was convenient for such testing due to easily replaceable wings and use of control surfaces mounted exclusively on the tail. The UAS was retrofitted with a supercritical wing aimed at achieving a flutter speed of Mach 0.95 at FL250, and used a custom control system for flutter suppression. \footref{fn:DAST} To this day, the DAST aircraft is the only unmanned aircraft that approaches the proposed flight condition of HUGO in terms of flight speed and mission. \\
The X-56A is a NASA testbed flying-wing UAS used to test aeroelastic phenomena on high aspect ratio wings in a low-speed flight regime, with a takeoff speed of 65 kts (33 m/s), maximum airspeed of 135 kts (69 m/s) during cruise and 150 kts (77 m/s) in a dive \cite{Schaefer2023}. Although the aircraft did not approach the intended flight cruise speed of HUGO, its mass and size are closer to that intended for HUGO.\\
Similarly, the FLEXOP UAS was used as a modular testbed with interchangeable wings for researching aeroelastic phenomena and design of flutter-suppressing control systems. \cite{Bauer2020} \cite{Yu2022} Analogous to the NASA X-56A aircraft, the cruise speed of $38 m/s$ of FLEXOP is significantly below the cruise speed of HUGO. However, the mission and scale of this aircraft approach HUGO much more closely, making it a reasonable reference for initial sizing. Furthermore, the scalability of the high wing-loading structure of DAST is uncertain and may not be within realistic design bounds for HUGO. Therefore, the X-56 and FLEXOP aircraft are also included in the reference aircraft list.

The design lift coefficient $C_{L_\text{des}}$ is determined based on an initial wing loading $W/S$ estimate and requirements from \cite{BaselineReport} at the sepcified cruise flight condition. Parameters of reference aircraft, including the wing loading itself, landing speed and cruise conditions are used in the estimate.\\
HUGO has to perform high subsonic flight and take-off or land on a paved runway. The initial estimate of $W/S = 1000\  N \cdot m^{-2}$ allows for sufficiently low stall speed while remaining a reasonable design value by lying between the considered reference aircraft. The take-off and landing of the X-56A and FLEXOP were performed from on paved runway. The DAST drone was launched from a carrier aircraft and captured by helicopter after completing a test flight. This allowed for a smaller wing due to the absence of ground takeoff and landing. \footref{fn:DAST} The high wing loading achieved in the DAST UAS sets an upper bound for the HUGO wing design. A significantly smaller wing is needed for HUGO compared to the X-56A or FLEXOP aircraft due to the high-speed cruise condition, yet the wing still must allow for take-off and landing at a reasonable speed. The mass $m$ and aspect ratio $AR$ are determined based on the mass requirement \textbf{HUG-GEN-001} and the aspect ratio requirement \textbf{HUG-AER-002} \cite{BaselineReport}.  The average value $AR=\frac{27+19}{2}=23$ from the range given in \textbf{HUG-AER-002} was used. The wing area $S$ can then be obtained from the wing loading. From these parameters, the wing span $b$ and mean chord length $MAC$ can be derived. The taper ratio $\lambda = 0.5$ for HUGO was selected based on guidlines by Torenbeek \cite{Torenbeek1982} and lies within the upper bound considered for jet transport aircraft in a similar flight regime. A slightly larger taper ratio is introduced due to the high aspect ratio of the wing. Torenbeek suggests that the sweep angle is determined primarily by compressibility effects and should increase progressively above a Mach number of 0.8. \cite{Torenbeek1982} An initial sweep $\Lambda_{LE} = 15 \degree$ was selected, similar to that of the X-56A and FLEXOP aircraft. This is considered a good balance between obtaining reasonable low-speed lift characteristics and delaying drag divergence at the test velocity. Relevant reference aircraft parameters \footnote{\href{https://ntrs.nasa.gov/api/citations/20120015944/downloads/20120015944.pdf}{NASA X-56A parameters} (accesed on 8.5.2026)} and preliminary estimates for HUGO are shown in \autoref{tab:DAST_HUGO}. Several additional wing parameters were obtained using the AeroToolbox website \footnote{\href{https://aerotoolbox.com/wing-plot-tool/}{AeroToolbox wing plotting tool}(Accessed on 12.5.2026)} and \cite{Torenbeek1982}.

The ARW-2 wing loading is several orders of magnitude larger than that of the X-56A UAS. This is due to the launch and landing conditions of both aircraft.  After matching the required cruise and takeoff/landing conditions, an initial wing loading $W/S = 1000 \ N \cdot m^{-2}$ was selected for HUGO. This value reflects the cruise Mach number of 0.75 as well as the landing and takeoff conditions for HUGO, requiring a larger wing. At the same time, the DAST ARW-2 wing design proves this magnitude of wing loading can be designed and used. The initial wing planform is likely to be revisited at a later design iteration and serves as an initial design point.

\begin{table}[H]
    \centering
    \caption{Reference aircraft and HUGO initial parameters \cite{perry1982methodology} \cite{Eckstrom1994} \cite{Bauer2020} \cite{Yu2022}} 
    \label{tab:DAST_HUGO}
    
    \begin{tabular}{|c|c||c|c|c|c|}
    \hline
        \textbf{Parameter} & \textbf{Symbol} & \textbf{DAST ARW-2} & \textbf{X-56A} & \textbf{FLEXOP} &\textbf{HUGO} \\
        \hline\hline
        Wing span [m] & $b$ & 5.79 & 8.53 & 7 & 3.36 \\
        \hline
        MAC [m] & $\overline{c}$ & 0.56 & 0.61 & 0.35 & 0.15 \\ 
        \hline
        Aspect ratio [$-$] & $AR$ & 10.3 & 14 & 20 & 23 \\
        \hline
        Taper ratio [$-$] & $\lambda$ & 0.4 & 0.46 & 0.5 & 0.5\\
        \hline
        Leading edge sweep [$\degree$] & $\Lambda_{LE}$ & 28.8 & 22 & 20 & 15\\
        \hline
        Wing area [m$^2$] & $S$ & 3.25 & 5.20 & 2.45 & 0.49 \\
        \hline
        Mass (maximum take-off) [kg] & $m$ & 1060 & 238 & 65 & 50 \\
        \hline
        Wing loading [N$\cdot$m$^{-2}$] & $W/S$ & 3216 & 449 & 260 & 1000 \\
        \hline
    \end{tabular}
\end{table}



\subsection{Design lift coefficient}
\label{subsec:cl_des}
To estimate $C_{L_\text{des}}$, \autoref{eq:C_L} can be applied for a cruise air density $\rho_{cr}$ of $0.5112 \ kg \cdot m^{-3}$ at FL270, and a cruise speed $V_{cr}$ of $230.2945 \ m \cdot s^{-1}$ corresponding to the test condition. Using the initial value for $W/S$ from \autoref{tab:DAST_HUGO}, an initial  $C_{L_\text{des}}$ of 0.074 was obtained.
%It includes the air density at the altitude of 27,000 feet $\rho_{cr}$ and the airspeed $V_{cr}$, for which it holds that $V_{cr} = M \cdot \sqrt{\gamma R T}$ at the ISA altitude of 27,000 ft.

\begin{equation}
    C_{L_\text{des}} = \frac{\left(\frac{W}{S}\right)}{\frac{1}{2} \rho_{cr} V_{cr}^2} \approx 0.074
    \label{eq:C_L}
\end{equation}

\subsection{Cruise Reynolds number}
Based on the initial estimates from \autoref{tab:DAST_HUGO}, the Reynolds number for cruise $Re_{cr}$ and stall condition $Re_\text{stall}$ was calculated using \autoref{eq:Re}.
\begin{equation}
    Re = \frac{\rho \overline{c}V}{\mu} \approx 1.15\cdot10^6
    \label{eq:Re}
\end{equation}
Using the cruise conditions $\rho_{cr}$, $V_{cr}$ and dynamic viscosity $\mu_{cr}$ corresponding to Torenbeek's ISA model \cite{AdvancedACDesign_Torenbeek}, a cruise Reynold's number $Re_{cr}$ of $1.15\cdot10^6$ was obtained for test conditions at FL270. The MAC $\overline{c} \approx 0.15 \ m$ is calculated from the wing parameters in \autoref{tab:DAST_HUGO}.

\subsection{Maximum lift coefficient}
The maximum lift coefficient, $C_{L_\text{max}}$ has to be determined to obtain the stall speed of HUGO, $V_\text{stall}$. As stated previously, the use of a supercritical wing profile is highly recommended for the considered test conditions. A preliminary estimate of  $C_{L_\text{max}}$ is made by considering a NASA SC(2)-0010 airfoil and the DATCOM method for estimating the maximum lift coefficient for finite wings. 

An estimation of $C_{L_\text{max}}$ can be made using Xfoil software\footnote{\href{http://airfoiltools.com/polar/details?polar=xf-sc20412-il-500000-n5}{AirfoilTools}(Accessed on 07.05.2026)} and an expression for estimating the wing $C_{L_{max}}$ from the $C_{l_{max}}$ of the wing profile. For a $Re \approx 1.15\cdot10^6$ and a critical amplification factor $N$ of 5, a $C_{l_\text{max}} \approx 1.4$ was found \footnote{\href{http://www.airfoiltools.com/polar/details?polar=xf-sc20010-il-1000000-n5}{NASA SC(2)-0010 airfoil}(Accessed on 11.5.2026)}. $N$ is deemed suitable for this usecase due to likely disturbed airflow from the wing surface roughness. \cite{Koul} Then, according to Torenbeek \cite{Torenbeek1982}, \autoref{eq:CL_max} is used to obtain $C_{L_{max}}$. \autoref{eq:Lambda_conversion} is used to convert from the leading edge sweep angle $\Lambda_{LE}$ to the quarter-chord sweep angle $\Lambda_{0.25c}$ and derived from the expression in \cite{Torenbeek1982}.

\begin{equation}
    \Lambda_{0.25}
    =
    \tan^{-1}
    \left(
    \tan \Lambda_{\mathrm{LE}}
    -
    \frac{1-\lambda}{AR (1+\lambda)}
    \right)
    \label{eq:Lambda_conversion}
\end{equation}

\begin{equation}
    C_{L_{\max}} = 0.9 \cdot C_{l_{\max}} \cdot \cos\Lambda_{0.25c}
    \label{eq:CL_max}
\end{equation}

This leads to $\Lambda_{0.25c} \approx 14.22$ and $C_{L_{\max}} \approx 1.22$ without high-lift devices.\\ \todo{cite this stuff}
\todo[inline]{In the trade-off chapter, we use $\lambda$ for eigenvalues, and here again for something else. Perhaps swapping one of them would be a good idea.}
\subsection{Stall speed}
The next step is to identify the preliminary stall speed of HUGO. With a maximum lift coefficient $C_{L_{max}} \approx 1.22$ and $S = 0.49 \ m^2$ at a mass $m = 50 \ kg$ and sea level ISA conditions, the stall speed can be calculated using \autoref{eq:Vstall}.
\begin{equation}
    V_\text{stall} = \sqrt{\frac{2W}{\rho S C_{L_\text{max}}}}
    \label{eq:Vstall}
\end{equation}
From \autoref{eq:Vstall}, a stall speed $V_{stall} \approx 37 \ m \cdot s^{-1}$. Comparing this to the stall speed of the X-56A UAS of 33 m/s, the result appears reasonable for an initial estimation.\\

\subsection{Supercritical airfoil}
Since the aircraft is expected to operate at Mach 0.75, it is necessary to take transonic phenomena into account. When flying in this region with a traditional airfoil, the airflow might locally be supersonic, which creates strong shock waves on the wing \cite{Whitcomb1974}. This leads to a large increase in drag, which is detrimental for the performance of the aircraft and introduces boundary layer separation. Such effects can be mitigated by the use of a supercritical airfoil. 

A supercritical airfoil is capable to operate at transonic conditions with only the onset of weak shock waves \cite{Harris1990}. An example of such airfoils used in the preliminary design for HUGO is a supercritical SC(2)-XXYY family developed by NASA, where XX is the airfoil design lift coefficient, and YY is the maximum thickness in percent \cite{Harris1990}. The advantage of using this system is the inclusion of the design lift coefficient, which makes it possible to narrow down the choices for HUGO using $C_{L_\text{des}}$ values found in \autoref{subsec:cl_des}.

\subsection{Airfoil selection}
Based on the data obtained in the previous subsections, an airfoil from the NASA database was selected. The ARW-2 wing of the DAST drone used supercritical airfoils with varying thicknesses along the wing span. The profile thickness varied between $\frac{t}{c} = 0.14$ at the wing root to $\frac{t}{c} = 0.10$ at the wing tip. \cite{ARW-2_NASA} Based on the initial design lift coefficient $C_{L_\text{des}} = 0.074$, the selected airfoil is based on a NASA report on supercritical airfoils. \cite{Harris1990} It was decided that the the NASA SC(2)-0012 will be used as the initial wing profile. The selected airfoil is symmetrical and designed for supercritical conditions.\\
A symmetrical supercritical airfoil is beneficial for low lift coefficient flight conditions. Only airfoils designed for $C_{l_{des}} = 0.4$ were tested, which is already significantly above the value suitable for HUGO. The symmetrical airfoil generates the desired lift coefficient with a low angle of attack. Furthermore, significant negative pitching moment is generated by non-symmetrical supercritical airfoils, like the NASA SC(2)-0412 \footnote{\href{http://airfoiltools.com/polar/details?polar=xf-sc20412-il-1000000-n5}{NASA SC(2)-0412}(Accessed on 11.5.2026)}, which could result in a larger and, consequently, heavier empannage. The reason for selecting a high thickness-to-chord ratio $t/c$ is the expected bending load on the slender wing of HUGO. The increased area moment of inertia of a thicker wing generally helps bending resistance. There are several risks associated with these aifoils with a relative high $\frac{t}{c}$ ratio. Compared to airfoils with higher $C_{L_\text{des}}$, flow separation and weak shock waves appear earlier, further from the trailing edge. However, at the same time, it has been reported that the overall drag divergence Mach number for airfoils with lower $C_{L_\text{des}}$ is generally higher. For the selected airfoil, the drag divergence Mach number $M_{dd}$ has been found to be around $M_{dd} = 0.8$, which is beyond the conditions HUGO will operate. \cite{Harris1990}

\section{Aerodynamic modelling}
In order to perform preliminary sizing, aerodynamic analyses of the aircraft were performed by means of the lifting line theory (LLT) with modifications for nonzero sweep provided by the AeroSandbox Python package \footnote{\href{https://aerosandbox.readthedocs.io/en/master/autoapi/aerosandbox/aerodynamics/aero_3D/lifting_line/index.html}{AeroSandbox} (accessed 18.05.2026)}. The goals of these analyses were to obtain the lift and induced drag of the aircraft, required for the later calculation of the $\frac{L}{D}$ ratio, the lift curve slopes of all the lifting surfaces (both individually and combined with other lifting surfaces), the lift distribution along the span, the calculation of the aerodynamic centre position, and the determination of the angle of attack required for trimming. This section first outlines the generation of the model, followed by the description of the performed analyses and their verification.

\textcolor{red}{Marek's inviscid modifications of the AeroSandbox source code}. The inviscid analysis has been verified by comparing its output with that of the standard viscous analysis. That is, by imposing the maximum relative difference of less than 1\% in the lift coefficient.

To begin with, the model is created parametrically, that is, as a function of the input parameters. These include the planform parameters of lifting surfaces, the chosen set of airfoils, and relative positions and orientations between the lifting surfaces. A lifting surface is generated by specifying the properties of an arbitrary number of sections along its span. Each section is, in turn, specified by the position of its leading edge, chord length, the airfoil, and twist. Specifically, the airfoil used in the aforementioned analyses was the same for all sections of all lifting surfaces, namely NASA SC(2)-0012. The measurement data from \cite{Mineck1987} has been used to obtain the lift curve slope of the airfoil, namely $
0.125\frac{1}{\degree}$  The sections in this model are spaced uniformly to simplify the extraction of the lift distribution results. Once all lifting surfaces are specified, the aircraft model is generated by combining the surfaces in a single airplane object. This object can include a single lifting surface (the main wing in isolation, for example), as well as all the lifting surfaces combined (main wing, canards, empennage).

The airplane model is subsequently used as an input to the LLT-based analyses. In addition, atmospheric properties (air temperature and density) and flight conditions (angle of attack, sideslip angle, airspeed and rotational rates) are required as inputs. Since the analysis results do not include the lift distribution, the position of vortex centres and vortex strengths is extracted from the source code. Then, the Kutta–Joukowski theorem ($L=\rho V\Gamma$) is used to compute lift per unit span at each section. In order to find the lift curve slope and the aerodynamic centre position of a lifting surface, the analysis is run for a range of angles of attack (two being sufficient but more of them actually used to rule out inconsistencies). The lift curve slope is then simply the slope of the obtained $C_L-\alpha$ pairs, whereas the aerodynamic centre is the point about which the aerodynamic moment is constant (calculated via \autoref{eq:aerodynamic_centre}).

\begin{equation}
    x_{ac} = x_{ref}-c_{ref}\cdot \frac{dC_m}{d_{\alpha}}
    \label{eq:aerodynamic_centre}
\end{equation}

Furthermore, a more difficult analysis problem is to run the trim analysis, since it requires a specific angle of attack. While the definition of the trim condition entails both vertical force and pitch moment equilibria, this analysis only considers vertical force equilibrium due to the preliminary nature of the design (no elevators, empennage assumed fixed for simplicity). Therefore, in order to find the required angle of attack, two approaches have been employed. In the first approach, the angle of attack is calculated by finding the lift force required for a specific flight condition, running the simulation for a range (at least two) angles of attack, and then calculating the angle of for obtaining the required lift, while assuming no lift at $\alpha=0\deg$. The second approach uses the Aerosandbox's optimizer with the objective of finding the angle of attack at which the obtained total lift equals the weight. Nevertheless, the second approach is not only more computationally expensive, but also more difficult to verify and validate. Therefore, only the first approach will be employed.

In order to verify the results, unit tests have been written (as can be found in the repository). Some of them are pure code verification, so these will not be discussed here. 

\textcolor{red}{Here table with configurations}

The slope of the lift curve of a high-aspect-ratio wing in the incompressible regime has been compared with Prandtl's formula for high-aspect-ratio wings that assumes elliptical lift distribution (\autoref{eq:prandtl_slope}). The error for this case is about 19.1\%, which can be explained by two things. Firstly, Prandtl's formula assumes elliptical lift distribution, which is not the case for this rectangular wing. Secondly, the LLT assumes a thin airfoil, which is not the case here. In fact, the supercritical airfoil used here has a relatively large thickness-to-chord ratio 

\begin{equation}
    a=\frac{a_0}{1+\frac{a_0}{\pi AR}}
    \label{eq:prandtl_slope}
\end{equation}

Furthermore, the lift curve slope of the second wing model in high subsonic compressible flow is compared with that of the DATCOM method \footnote{\href{https://www.fzt.haw-hamburg.de/pers/Scholz/HOOU/AircraftDesign_7_WingDesign.pdf}{Scholz} (accessed 18.05.2026)} \autoref{eq:datcom_slope}. The result is much better than for the Prandtl formula, with the relative difference between the results being 10.5\%. That implies that the Prandtl formula overly simplifies the aerodynamics.

\begin{equation}
    a
=
\frac{2\pi AR}
{2 + \sqrt{
4 + \left(\frac{AR\beta}{\eta}\right)^2
\left(1 + \frac{\tan^2 \Lambda_{0.5C}}{\beta^2}\right)
}}
\qquad
\beta = \sqrt{1 - M^2}
\label{eq:datcom_slope}
\end{equation}

The analysis computing the AoA required for trim is verified in a similar way...
\textcolor{red}{Possibly add plots here}
+verification aerodynamic centre, aoa sweep


\section{Empennage preliminary design}
The next step is a preliminary sizing of the empennage. Since the NASA X-56A is a tailless aircraft, it is not included in the list of reference aircraft values. Instead, the Boeing truss-braced wing aircraft concept is used. \cite{XiongBoeing} First, data about the horizontal and vertical tail from several reference aircraft were collected, shown in \autoref{tab:tail_params}. Several tail parameters were not available directly. Therefore, drawings and projected views have been used to record tail parameters using known values, such as the wing span. Note that HT refers to the horizontal tailplane and VT refers to the vertical tailplane.

\begin{table}[H]
    \centering
    \caption{Tail parameters of reference aircraft and HUGO \cite{XiongBoeing} \cite{Yu2022} \cite{Bauer2020} \cite{perry1982methodology} \cite{Eckstrom1994}} 
    \label{tab:tail_params}
    
    \begin{tabular}{|c|c||c|c|c|c|}
    \hline
    \textbf{Parameter} & \textbf{Symbol} & \textbf{DAST ARW-2} & \textbf{FLEXOP} & \textbf{TTBW} & \textbf{HUGO} \\
    \hline\hline
    Tail configuration [-] & - & Conventional & V-tail & T-Tail & T-Tail\\
    \hline
    HT span [m] & $b_h$ & 1.71 & 1.27 (projected) & 11.75 &\\
    \hline
    HT arm [m] & $l_h$ & 2.66 & 2.0 & 21.09 &\\
    \hline
    Non-dimensionalized HT arm [-] & $\frac{l_h}{\overline{c}}$ & 4.75 & 2.67 & 7.25& \\
    \hline
    HT aspect ratio [-] & $AR_h$ & 3.5 & 4.2 & 4.52 &\\
    \hline
    HT taper ratio [-] & $\lambda_h$ & 0.4 & 0.52 & 0.35 &\\
    \hline
    HT dihedral [\degree] & $\Gamma_h$ & N/A & 35 & N/A &\\
    \hline
    HT leading edge sweep [\degree] & $\Lambda_{LE_{h}}$ & 45 & 23.55 & 21.3&\\
    \hline
    HT quarter chord sweep [\degree] & $\Lambda_{qc_{h}}$ & 41.27 & 19.83 & 15.82 &\\
    \hline
    HT area $[m^2]$ & $S_h$ & 0.836 & 0.39 & 30.54 & \\
    \hline
    VT half-span [m] & $b_v$ & 0.83 & N/A & 9.64 &\\
    \hline
    VT leading edge sweep [\degree] & $\Lambda_{LE_{v}}$ & 52 & N/A & 16.6 & \\
    \hline
    \end{tabular}
\end{table}

Since the aircraft has to be able to allow for stable, neutrally stable, and unstable flight configurations, the neutral position of the c.g. is chosen as the design point. However, the aircraft has to be controllable at the same time, with a certain margin. These two conflicting properties of stability and controllability are governed by the equations shown in \autoref{eq:Stability} and \autoref{eq:Controllability}, respectively. Clearly, they describe the positions of the centre of gravity as a function of the aerodynamic centre position and numerous aerodynamic and geometric parameters. The two curves constitute the so-called scissor plot, where the aircraft is controllable when ${\bar{x}_{cg}}>{\bar{x}_{cg}}_{\text{controllability}} $ and stable when ${\bar{x}_{cg}}<{\bar{x}_{cg}}_{\text{stability}}$. Hence, in order to constrain the design space, we enforce $\bar{x}_{cg}={\bar{x}_{cg}}_{\text{stability}}$ and $\bar{x}_{cg}-{\bar{x}_{cg}}_{\text{controllability}}=\text{C.G.E.}$, where  C.G.E. is the minimum required excursion between the controllability curve and the design point. 

Also, a certain range of the c.g. positions is required. All of that is tied to the position and size of the horizontal stabilizer through the so-called scissor plot. The canard configuration is not considered here, since not considering it makes the design conservative. 

\begin{equation}
{\bar{x}_{cg}}_{\text{controllability}} = \bar{x}_{ac} - \frac{C_{m_{ac}}}{C_{L_{A-h}}} + \frac{C_{L_h}}{C_{L_{A-h}}} \frac{S_h l_h}{S\bar{c}} \left(\frac{V_h}{V}\right)^2
\label{eq:Controllability}
\end{equation}

\begin{equation}
{\bar{x}_{cg}}_{\text{stability}} = \bar{x}_{ac} + \frac{C_{L_{\alpha_h}}}{C_{L_{\alpha_{A-h}}}} \left(1 - \frac{d\varepsilon}{d\alpha}\right) \frac{S_h l_h}{S\bar{c}} \left(\frac{V_h}{V}\right)^2
\label{eq:Stability}
\end{equation}

\section{Landing Gear Configuration}
\label{sec:Landing_gear_config}

This section presents the selection analysis for the possible landing gear configurations of HUGO. Out of the five configurations, two of them use a buried design, with the wheels retracting vertically, while the other three configurations have a sideways diagonal retracting platform. 

The two undercarriage concepts differ in both retraction mechanism and stowage arrangement. In the external side-pod configuration, the landing gear is housed in nacelles mounted outside the fuselage. This arrangement can reduce the intrusion of the undercarriage into the main fuselage structure and can simplify the internal layout. However, the addition of side pods increases wetted area and introduces interference drag at the pod fuselage junction.

In the internal fuselage-bay configuration, the landing gear is retracted into the fuselage structure. This provides a cleaner external geometry, but requires dedicated internal volume, local structural reinforcement, landing gear doors, and a more integrated retraction mechanism. The placement of the undercarriage is also constrained by the need to avoid interference with the wing box, especially if a modular wing structure is used. Therefore, the final landing gear arrangement will depend strongly on the selected HUGO configuration and on the available structural integration space.

%Since all five aircraft configurations use the same general landing gear concept, namely a retractable undercarriage, it does not directly affect the main trade-off selection.

In addition to the design requirements determined earlier, HUGO must also meet ground stability and manoeuvrability requirements. \cite{roskamLandingGearSystems} and \cite{Vos2025AirplaneDesignAnalysis} state that for a landing gear to be operational, the following constraints must be obeyed:
\begin{itemize}
    \item The scraping angle must be equal, to or larger then $15 \deg$ (pitch angle req.)
    \item The angle between the vertical line and the line through the main gear and the C.G. has to be at least equal to the scraping angle (pitch angle req.)
    \item The turn over angle must be less than 55 degrees;
    \item The wing-tip-to-landing-gear-wheel angle must be greater than 7 degrees
    \item The front landing gear must bear between 8\% and 15\% of the total aircraft weight
\end{itemize}
\todo{roskam picture}

By using trigonometric relations between the landing gear and critical points such as the location of the C.G., the positions of the beginning and end of the tail cone, and the wing tips as described in \cite{roskamLandingGearSystems}, the location of the main and nose landing gear, as well as the landing gear length, and main landing gear separation can be optimised for the given geometric parameters of the fuselage. Using Sequential Least Squares Quadratic Programming (SLSQP), the strut length, x-position, and the lateral track of the main landing gear, as well as the x-position of the nose landing gear can be minimised. 


Pros and cons for the 2 LG 
1. Outside nacelle: 

2. Internal bay:


%The landing gear has to comply with a couple of conditions in order to ensure proper functionality. The tip-back angle of HUGO is $\theta =15\deg$\cite{roskamLandingGearSystems} and a requirement from \cite{roskamLandingGearSystems} states that the angle between the vertical line and the line through the main gear and the C.G. has to be larger 



\section{Structural Analysis}

In order to perform the trade-off analysis, an initial structural analysis of loaded components must be performed in order to estimate the weight penalty of versatility for each configuration. The components under significant load mainly include the fuselage, main wing, canard, and empennage. 

\begin{math}
    y_{c1} = R \left[ 1 - \frac{2 \sin \alpha}{3 \alpha} \left( 1 - \frac{t}{R} + \frac{1}{2 - t/R} \right) \right]
\end{math}

\begin{math}
    y_{c} = R - y_{c1}
\end{math}

\begin{math}
    y_{c} = R * \frac{2 \sin \alpha}{3 \alpha} \left( 1 - \frac{t}{R} + \frac{1}{2 - t/R} \right) 
\end{math}

\begin{math}
I_x = R^3 t \left[ \left( 1 - \frac{3t}{2R} + \frac{t^2}{R^2} - \frac{t^3}{4R^3} \right) \times \left( \alpha + \sin \alpha \cos \alpha - \frac{2 \sin^2 \alpha}{\alpha} \right) + \frac{t^2 \sin^2 \alpha}{3R^2 \alpha (2 - t/R)} \left( 1 - \frac{t}{R} + \frac{t^2}{6R^2} \right) \right]
\end{math}

\section{Material Selection}
This Section addresses the selection of aerospace materials for various subsystems of HUGO. Key factors considered in this analysis are: structural properties, material cost, manufacturing cost, durability, sustainability, and special properties. For the purpose of a preliminary material selection, the aircraft was divided into the following subgroups, all of which are expected to be made primarily from one material:

%The aircraft will most likely not be made of one material but instead a combination of several. The aircraft can be roughly be split up into the following: fuselage structure, wing structure, empennage and control surfaces, and landing gear. It is a simplification to assume that each of those subcategories would by made of the same material but this is a valid assumption at this stage of the design process. The specific operating requirements and conditions can be detailed for each subsystem. 

%\autoref{tab:aerospace_applications} lists the main loading modes and constraints for some of the major aircraft components. 

\begin{comment}
\begin{table}[H]
\centering
\caption{Aerospace Applications/Idealizations, adapted from \cite[p. 6] {Arnold2012MaterialsSelection}}
\begin{tabular}{|c|p{7cm}|p{6.5cm}|}
\hline
\textbf{Application} & \textbf{Loading} & \textbf{Constraints} \\ \hline

Truss & Tension & Strength \\ \hline

Fuselage  & Compression, Tension, Bending  & Buckling, Strength, Deflection \\ \hline

%Shafts & Thin walled tube/Beam & Torque, Rotation, Bending  & Bending, Shear, Vibration \\ \hline

%Combustion Chambers & Thin walled Cylinder, Internal pressure & Pressure, Noise & Thermal, Oxidation \\ \hline

%Blades (Fan, Compressor, Turbine) & Tie/Beam & Bending, Tension, Fatigue, Thermal, Creep & Deflection, Strength, Vibration, Erosion, Temperature \\ \hline

%Disks (Compressor, Turbine) & Disk & Rotation, Thermomechanical Fatigue & Strength, Ductility, Toughness, Oxidation \\ \hline

%Space Truss Framework & Tie rod & Axial Loading & Strength, Buckling \\ \hline

%Nozzle & Pressure, Noise, Thermal & Strength, Vibration \\ \hline

%Wing & Bending, Tension, Twist, Vibration, Fatigue & Buckling, Vibration, Strength, Deflection \\ \hline

%Cryogenic Tanks & Spherical Shell & Internal Pressure, Thermal & Strength, low temperature Yield before break, Leak before break \\ \hline

\end{tabular}
\label{tab:aerospace_applications}
\end{table}
\end{comment}

 
\begin{itemize}
    \item \textbf{Fuselage structure:}
    The fuselage is the main structure of the aircraft, which acts as a tethering point for the wings, empennage, and propulsion \cite{CHATTERJEE2019199}. Additionally, it contains the major control system and payload. The main loads acting on a fuselage are compression, tension, and bending, and main design constraints are based on phenomena such as buckling, yielding, and deflection \cite{Arnold2012MaterialsSelection}. For HUGO, the largest loads are expected to occur at a flight speed of Mach 0.75.
    
    \item \textbf{Wing structure:} The wings of an aircraft are subjected to varying types of loads during flight, including bending, tension, compression, twist, vibration, and fatigue \cite{Arnold2012MaterialsSelection}. The main design constraints are related to an analysis of buckling, vibrations, yielding, and deflection. The material for HUGO's wing planform shall therefore be chosen to provide a sufficiently stiff and strong structure (e.g. to withstand the load factors defined in the flight envelope shown in \autoref{fig:V-n_diagram}). Depending on the nature of the test, and especially the required stiffness, the choice of material for this subsystem is expected to vary.

    \item \textbf{Empennage and control surfaces:} The empennage will experience loads of varying degrees during the flight. The structure needs act as a moment-arm for the generated forces at the tailplane. This makes it possible to generate moments around the centre of gravity both in the vertical and horizontal plane. This is not only to ensure straight and steady flight but also to make manoeuvring (change of heading, flight direction and etc.) possible for the aircraft.   

    
    \item \textbf{Landing gear:} 
    The landing gear must be able to carry the weight of the aircraft, as well as withstand take-off and landing loads. According to Albuquerque\cite{structural_loads_handbook}, the key loads acting on the undercarriage can be split into: ground loads (static load, braking, braked roll, ground turns, reversed braking, nose wheel yaw and steering, pivoting, towing etc.), landing loads, and take-off loads (associated with retraction of the landing gear and gyroscopic effects of changing wheel direction). All of these cases are relevant for HUGO, as the in its final configuration, the landing gear retracts into the aircraft during flight (see: \autoref{sec:trade-off_final_design}).  

    \item \textbf{Pylons:} For the design of HUGO, it is assumed that an already existing engine will be chosen, together with a nacelle provided by the manufacturer. However, the design of pylons still has to be considered. The main loads acting on the pylons come from the engines, specifically the associated vibrations and thermal loads \cite{pylon_load_analysis}. Therefore, key considerations for this structural element are related to fatigue life and resistance to high temperatures.    

    \item \textbf{Wing/Empennage Ports:} The modular ports defined for the empennage and/or wing structures of HUGO will have to withstand the loads acting at the connection points between the fuselage and the given planform. Structural analysis shows that both the bending and shear loads are the highest at the root \cite{structural_loads_handbook}. Therefore, the material for the ports shall be chosen such that the resistance to the aforementioned load types is prioritised.
\end{itemize}

It is important to note that at this design stage, the materials are only chosen for the primary components of each structural subsystem. A more detailed design of elements such as antenna mounts, access panel hardpoints, fastener inserts etc., as well as the details of internal supporting structure, will be analysed in more detail once the final external configuration is established for HUGO. 


\subsection{Material Analysis of Reference Aircraft}
\label{sub:materials_reference_aircraft}
The material selection shall begin with an analysis of reference aircraft, X56-A and DASR ARW-2 previously introduced in \autoref{sub:lift_reference_aircraft}, as well as existing trends in the aerospace industry. The purpose of this research is to ensure that HUGO can model the behaviour of contemporary aircraft as accurately as possible.

X56-A, the high aspect ratio testbed developed by NASA in the 2010s, features a carbon fibre composite centrebody with honeycomb bulkheads \cite{X56-AMaterials}. Carbon fibre composites were also used for the wing spars and ribs, as well as the skin for stiff wings. Flexible wing configurations featured glass fibre skins due to their lower stiffness. It also features a stiff landing gear designed for parachute landing loads. While specific materials used for the landing gear do not seem to be publicly available, accessible footage suggests that the structure is metallic.

The DAST ARW-2, a slightly older concept from the 1980s, featured a wing using 7075-T73 aluminium spars and ribs covered with a fibreglass honeycomb sandwich skin material (36 layers at inboard end, 27 at outboard, with 25-percent of the layers at $\pm$45 degree orientation) \cite{Eckstrom1986ARW2}. High stiffness was achieved by incorporating a thick aluminium plate in the centre section of the wing, and strain gauge bridges were mounted on the spars and ribs.

The prevalence for composite usage in external structures is also present in modern commercial aircraft. Leading aircraft manufacturers, such as Boeing and Airbus, both incorporate carbon laminates into the airframes of their products as a lighter alternative to metallic alloys\footnote{\href{https://www.airbus.com/en/newsroom/news/2017-02-aviations-material-evolution}{Airbus - Aviation's material evolution} (accessed on 13.05.2026)}\footnote{\href{https://www.boeing.com/content/dam/boeing/boeingdotcom/features/innovation-quarterly/2024/IQ-2024-Q3-full.pdf}{Boeing - 787 Dreamliner Difference} (accessed on 13.05.2026)}. Additionally, Airbus A330neo nacelles are made with 60\% composites, together with a honeycomb structure for acoustics\footnote{\href{https://www.safran-group.com/products-services/airbus-a330neo-nacelles}{Airbus A330neo nacelles} (accessed on 13.05.2026)}, while Boeing 777x uses fine metallic alloys instead\footnote{\href{https://www.safran-group.com/products-services/boeing-777x-exhaust-system}{Boeing 777x exhaust system} (accessed 13.05.2026)}. In the recent years, NLR has also began experimenting with implementing composites into upper spars of pylons    \footnote{\href{https://www.toraytac.com/newsroom/story/HelI/Thermoplastic-Upper-Spar-for-an-Aircraft-Engine-Pylon}{Automated Production of Large and Thick Thermoplastic Composite Aircraft Parts} (accessed on 18.05.2026)}. Aluminium, steel, and titanium alloys are also among the most commonly used materials for typical landing gear systems \cite{typicallandinggear}, as well as engine pylons (specifically Ti and Al-Li alloys)\cite{TiAl_pylon_materials}. 

\autoref{tab:materials_airliner_breakdown} showcases the usage of aerospace materials for two exemplary commercial aircraft: Airbus A350 XWB, and B787 Dreamliner. These models were chosen due to data availability, and were deemed as representative of modern aircraft systems.

\begin{table}[H]
   \caption{Materials used in external structures of A350 XWB\cite{skinmaterials_1} and B787\cite{skinmaterials_2}}
    \centering
    \begin{tabular}{|c|p{6cm}|p{6cm}|}
    \hline
        \textbf{Material} & \textbf{A350 XWB} & \textbf{B787} \\
        \hline\hline
         Carbon laminate & Fuselage skin, wings, centre wing box, keel beam, tail cone, skin panels, frames/stringers/doublers, doors & Fuselage skin, ailerons, flaps \\\hline
         Fibreglass & $-$ & wing skin, wing-to-fuselage joints\\\hline
         Al alloys & Cockpit, wing leading edges, engine inlets & Wing leading edges, engine inlets\\\hline
         Metallic alloys & Pylons (Ti) & Pylons (Al/Steel/Ti)\\\hline
         Carbon sandwich & $-$ & Nacelles, elevators, rudders\\ \hline
    \end{tabular}
    \label{tab:materials_airliner_breakdown}
\end{table}

%787 Dreamliner: Carbon laminate for fuselage skin + ailerons + flaps, fibreglass for wing skin and wing joints, aluminium for leading edges + engine inlets, aluminium/steel/titanium for pylons, carbon sandwich for nacelles, elevators and rudders. \cite{skinmaterials_2}

%Airbus A350 XWB: CFRP for wings, centre wing box + keel beam, tail cone, skin panels, frames/stringers/doublers, doors \cite{skinmaterials_1} aluminium/al-li for cockpit, leading edge of the wing, nacelle inlet, titanium pylons






\subsection{Aerospace Materials for Airframe}
\label{sub:materials_candidates}

According to Mouritz\cite{Mouritz2012}, there are no more than one hundred types of metal alloys, composites, polymers and ceramics that exhibit properties suitable for aerospace applications. The main materials used for airframes and engines include: aluminium alloys, titanium alloys, steels, and composites. Suitable representatives of each material group selected for the material design trade-off for HUGO are described in \autoref{tab:material_candidates_desc}.

\todo[inline]{Materials table with rationale and specific alloys/types}
\begin{table}[H]
    \centering
        \caption{Caption}
    \begin{tabular}{|c|c|p{3cm}|p{8cm}|}
    \hline
       \textbf{Material ID} & \textbf{Material} & \textbf{Alloy/Type} & \textbf{Description}\\\hline\hline
        HUGO-MAT-01 & Wood & Sitka spruce &  Wood has played a part in aerospace engineering since the first powered flight. The softwood Sitka spruce was chosen based on its historical use in the construction of aircraft stemming from high specific strength and stiffness \cite{bookJohn}. The material is additionally characterised by its low density and low cost. \\\hline
        HUGO-MAT-02 & Titanium & Ti-6Al-4V & \\\hline
        HUGO-MAT-03 & Aluminium & Al 7075-T76 &  Aluminium has been widely used in aviation since the 1930s, and has proved usefulness in aerospace industry. The specific alloy was chosen due to its outstanding anti-corrosive properties and stress corrosion cracking resistance \cite{LI2023944}. \\\hline
        HUGO-MAT-04 & Magnesium & Mg WE46-T6 & \\\hline
        HUGO-MAT-05 & Steel & Steel 18Ni (350)\tablefootnote{\href{https://www.matweb.com/search/DataSheet.aspx?MatGUID=d0e2e9edae8d4060aa22dc11dd1b66b5&ckck=1}{MatWeb AISI Grade 18Ni (350) Maraging Steel} (accessed 08.05.2026)} \cite{ANTOLOVICH1980717} & \\\hline
        HUGO-MAT-06 & Composite & Carbon/epoxy (quasi-direct.) & \\\hline
        HUGO-MAT-07 & Composite & Glass/epoxy (quasi-direct.) \tablefootnote{\href{https://www.goodfellow.com/eu/resources/e-glass-epoxy-material-information/}{E-glass/Epoxy - Material Information} (accessed 12.05.2026)} \cite{Vantadori} & \\\hline
        
    \end{tabular}
    \label{tab:material_candidates_desc}
\end{table}

%Since most metals used in aerospace occurs as some sort of alloy, the most relevant alloys has been chosen to limit the list of candidates \cite{Mouritz2012}.

Below is the list of the different materials that will be considered.   
\begin{itemize}
    %\item \textbf{Aluminium:} The most widely used material for aerospace in the last fifty years. It is characterised by its ease of manufacturing, high specific strength and stiffness. The alloy used for the analysis is Al (7075 Al-T76) \cite{LI2023944}. 
    %\item \textbf{Wood:} Wood has played a part in aerospace engineering since the first powered flight. The softwood Sitka spruce was chosen based on its historical use in the construction of aircraft stemming from high specific strength and stiffness \cite{bookJohn}. The material is additionally characterised by its low density and low cost. 
    \item \textbf{Magnesium} These alloys were used in the aerospace industry for a wide range of applications. It has the lowest density of all metals commonly used in aircraft. The alloy Mg (WE46-T6) will be used in the trade-off.
    \item \textbf{Titanium} It has a high strength at medium density. The $\beta$-Ti alloy (Ti-6Al-4 V) will be considered.  \cite{Pollock2016}. 
    \item \textbf{Steel} This has high strength with high density. The alloy used for the comparisons are Steel (18Ni (350))
    \item \textbf{Composites} This consists typically of two materials the matrix and the fibres. In the case of this the chosen composite the matrix is thermoset epoxy and the fibres are carbon or glass.    
\end{itemize}

A qualitative comparison of the candidate materials can be performed using relevant aerospace literature. \autoref{tab:materials_grading_properties} lists the grading of aerospace materials on key design factors, distinguishing between: very low (VL), low (L), medium (M), high (H) and very high (VH), based on Mouritz's classification of metal alloys and CFRP \cite[p. 6]{Mouritz2012}. Additional columns were added for wood and glass/epoxy composites based on a comparison to CFRP\footnote{\label{url:cfrp_vs_gfrp}\href{https://njnir.com/aerospace/cfrp-vs-gfrp?utm_source=chatgpt.com}{CFRP vs. GFRP in Aerospace Engineering: A Comparative Analysis of Composite Materials} (accessed on 19.05.2026)}, Mouritz material comparison \cite{Mouritz2012}, and material property plots made by Ashby et al \cite{ashby2019materials}. Based on the gathered data, a first-order comparison can be made between the considered material candidates.


 
\begin{table}[H]
\centering
\caption{Grading of aerospace materials on key design factors (adapted from \cite[p. 6]{Mouritz2012})}
\begin{tabular}{|l|c|c|c|c|c|c|c|}
\hline
\textbf{Property} & \textbf{Al} & \textbf{Ti} & \textbf{Mg} & \textbf{Steel} & \textbf{CFRP}  & \textbf{Glass/epoxy} & \textbf{Wood} \\
\hline\hline
Cost & L & H & M & M  & H & M & VL\\\hline
Weight (density) & L & M & VL & H  & VL & L & VL\\\hline
Stiffness (elastic modulus) & L/M & M & L & VH & H & L & L \\\hline
Strength (yield stress) & M & M/H & L & VH  & H & M & L \\\hline
Fracture toughness & M & H & L/M & L/M  & L & L & L\\\hline
Fatigue & L/M & H & L & M/H & H & M & H\\\hline
Corrosion resistance & M & H & L & L/M & VH & H & N/A \\\hline
High-temperature creep strength & L & M & L & H  & L & L & N/A\\\hline
Ease of recycling & H & M & M & H & VL & L & VH\\
\hline
\end{tabular}
\label{tab:materials_grading_properties}
\end{table}


\begin{comment}
\begin{table}[H]
\centering
\caption{Mechanical properties of selected aerospace materials (part of the table taken from chapter 17 table 17.1 in Introduction to Aerospace Materials \cite{Mouritz2012}).}
\label{tab:material-properties}
\renewcommand{\arraystretch}{1.2}
\begin{tabular}{|l|l|c|c|c|c|}
\hline
\textbf{ID}
  & \makecell{\\\textbf{Density}\\(g\,cm$^{-3}$)}
  & \makecell{\textbf{Specific}\\\textbf{modulus}$^*$\\$\times10^{6}$\\(N\,m\,kg$^{-1}$)}
  & \makecell{\textbf{Specific}\\\textbf{strength}\\$\times10^{3}$\\(N\,m\,kg$^{-1}$)}
  & \makecell{\textbf{Specific}\\\textbf{fracture}\\\textbf{toughness}\\$\times10^{3}$\\(N\,m$^{1.5}$\,kg$^{-1}$)} \\
\hline\hline
HUGO-MAT-01 &  0.4 & 22.5 & 138   & 10.0 \\ \hline
HUGO-MAT-02 & 4.6 & 23.5 & 217   & 21.7 \\ \hline
HUGO-MAT-03 &  2.7 & 25.9 & 174   & 10.7 \\ \hline
HUGO-MAT-04 & 1.7 & 26.5 & 117   & 11.8 \\ \hline
HUGO-MAT-05  & 8.1 & 24.6 & 295.6 & 4.9 \\ \hline
HUGO-MAT-06  & 1.7 & 29.4 & 450   & 22.4 \\ \hline
HUGO-MAT-07 & 2.0 & 28.7&  384  & 31.0 \\
\hline
\end{tabular}
\end{table}
\end{comment}

The base material properties for each of the materials considered are listed in \autoref{tab:base-material-properties}. However, they do not paint a full picture of the strengths and weaknesses of each material. In order to get a clearer picture, the specific modulus ($E/\rho$), specific strength ($\sigma_y/\rho$), specific fracture toughness ($K_{Ic}/\rho$), column buckling efficiency ($\sqrt{E}/\rho$), and panel/skin buckling efficiency ($\sqrt[3]{E}/\rho$) were calculated, and are listed in \autoref{tab:derived-material-properties}. These were then normalised and summed to arrive at an overall index for the mechanical properties of the materials. This was then used to rank the materials in the trade off. The details of these calculations can be found in this \href{https://docs.google.com/spreadsheets/d/1GpVjdeAKDGPKHCRgVgf5LDyppZcSMRnKjkgZY_FNGBo/edit?usp=sharing}{spreadsheet}. 

\begin{table}[H]
\centering
\caption{Base mechanical and economic properties of selected aerospace materials (Mechanical data derived from \cite{Mouritz2012}).}
\label{tab:base-material-properties}
\renewcommand{\arraystretch}{1.2}
\begin{tabular}{|l|c|c|c|c|}
\hline
\textbf{ID}
  & \makecell{\textbf{Density} ($\rho$)\\($kg\,m^{-3}$)}
  & \makecell{\textbf{Young's}\\\textbf{Modulus} ($E$)\\(GPa)}
  & \makecell{\textbf{Specific}\\\textbf{Cost}\\(USD\,$kg^{-1}$)}
  & \makecell{\textbf{Carbon}\\\textbf{Footprint}\\($CO_2\,kg^{-1}$)} \\
\hline\hline
HUGO-MAT-01\tablefootnote{\href{https://spotmaterial.com/material/spruce/sawnwood-kiln-dried/}{Spot material} (accessed 11.05.2026)} & 400  & 9000   & 0.85-1.7   & 0.05-0.5 \\ \hline
HUGO-MAT-02\tablefootnote{\href{https://spotmaterial.com/material/titanium/titanium-alloy-ti-6al-4v-grade-5-ti-6-4/}{Spot material} (accessed 11.05.2026)} & 4600 & 108100 & 21-25      & 8.1-26.5 \\ \hline
HUGO-MAT-03\tablefootnote{\href{https://spotmaterial.com/material/aluminium/7000-series-alznmgcu-7075-7475/}{Spot material} (accessed 11.05.2026)} & 2700 & 69930  & 2.55-5.1   & 4-20     \\ \hline
HUGO-MAT-04\tablefootnote{\href{https://spotmaterial.com/material/magnesium/cast-we-alloys-we43-we54/}{Spot material} (accessed 11.05.2026)} & 1850 & 49025  & 2.975-5.95 & 8.5-36   \\ \hline
HUGO-MAT-05\tablefootnote{\href{https://spotmaterial.com/material/alloy-steel/high-speed-tool-steel-hss-m-t/}{Spot material} (accessed 11.05.2026)} & 8100 & 199260 & 8-20       & 6.8      \\ \hline
HUGO-MAT-06\tablefootnote{\href{https://spotmaterial.com/material/carbon-fibre-cf/high-modulus-hm-ep-cf60-60-continuous-cf-hm/}{Spot material} (accessed 11.05.2026)} & 1700 & 49980  & 25.91-32.30& 17.5-23.2\\ \hline
HUGO-MAT-07\tablefootnote{\href{https://spotmaterial.com/material/high-modulus-glass-fibre/high-strength-s/}{Spot material} (accessed 12.05.2026)} & 2000 & 57400  & 6.5-8.7    & 4.8      \\
\hline
\end{tabular}
\end{table}

\begin{table}[H]
\centering
\caption{Derived specific material indices and economic factors for aerospace structural design.\tablefootnote{Full normalization calculations and trade-off data can be accessed via the project \href{https://docs.google.com/spreadsheets/d/1GpVjdeAKDGPKHCRgVgf5LDyppZcSMRnKjkgZY_FNGBo/edit?usp=sharing}{Spreadsheet}.}}
\label{tab:derived-material-properties}
\renewcommand{\arraystretch}{1.2}
\begin{tabular}{|l|c|c|c|c|c|c|}
\hline
\textbf{ID}
  & \makecell{\textbf{Specific}\\\textbf{Modulus}\\($E/\rho$)\\(GPa / $kg\,m^{-3}$)}
  & \makecell{\textbf{Column}\\\textbf{Buckling}\\($\sqrt{E}/\rho$)}
  & \makecell{\textbf{Plate}\\\textbf{Buckling}\\($\sqrt[3]{E}/\rho$)}
  & \makecell{\textbf{Specific}\\\textbf{Strength}\\($\sigma/\rho$)\\$\times10^{3}$\\(N\,m\,kg$^{-1}$)}
  & \makecell{\textbf{Specific}\\\textbf{Fracture}\\\textbf{Toughness}\\($K_{Ic}/\rho$)\\$\times10^{3}$\\(N\,m$^{1.5}$\,kg$^{-1}$)} 
  & \makecell{\textbf{Total}\\\textbf{Normalized}\\\textbf{Economic}\\\textbf{Factor}} \\
\hline\hline
HUGO-MAT-01 & 22.5 & 0.24 & 0.05 & 138   & 10.0 & 0.21 \\ \hline
HUGO-MAT-02 & 23.5 & 0.07 & 0.01 & 217   & 21.7 & 4.34 \\ \hline
HUGO-MAT-03 & 25.9 & 0.10 & 0.02 & 174   & 10.7 & 0.68 \\ \hline
HUGO-MAT-04 & 26.5 & 0.12 & 0.02 & 117   & 11.8 & 0.92 \\ \hline
HUGO-MAT-05 & 24.6 & 0.06 & 0.01 & 295.6 & 4.9  & 3.80 \\ \hline
HUGO-MAT-06 & 29.4 & 0.13 & 0.02 & 450   & 22.4 & 3.35 \\ \hline
HUGO-MAT-07 & 28.7 & 0.12 & 0.02 & 384   & 31.0 & 0.94 \\
\hline
\end{tabular}
\end{table}

\subsection{Economic factors}
\label{sub:materials_economics_sustainability}
The carbon footprints of each material, in Tonnes of CO2 per kg of material are listed in \autoref{tab:base-material-properties}. Cost is another feature to consider for the trade off. However, cost per kg does not paint a clear picture, since the differences in mechanical properties means that 1 kg of carbon fibre composite is not equivalent in strength to 1 kg of steel, for instance. To account for this, the cost per kilogram of each material was divided by its material indices, then normalized to a dimensionless value, and summed to arrive at the normalized material cost, where lower is better. The results of this can be seen in \autoref{tab:derived-material-properties}. The results of this analysis show that Sitka Spruce (HUGO-MAT-01) is the cheapest material in terms of raw cost, whereas Titanium (HUGO-MAT-02) is the most expensive. More details on this analysis are listed in the following \href{https://docs.google.com/spreadsheets/d/1GpVjdeAKDGPKHCRgVgf5LDyppZcSMRnKjkgZY_FNGBo/edit?usp=sharing}{spreadsheet}.



\begin{comment}

\begin{table}[H]
\centering
\caption{Pricing and emissions attributed to the selected aerospace materials}
\label{tab:material-economics}
\renewcommand{\arraystretch}{1.2}
\begin{tabular}{|l||c|c|}
\hline
\textbf{ID}
  & \makecell{\textbf{Specific}\\\textbf{cost}\\(euro\,$kg^{-1}$)}
  & \makecell{\textbf{Carbon}\\\textbf{Footprint}\\ ($CO_2 kg^{-1}$)} \\
\hline\hline
HUGO-MAT-01\tablefootnote{\href{https://spotmaterial.com/material/spruce/sawnwood-kiln-dried/}{Spot material} (accessed 11.05.2026)}                  & 0.85-1.7 & 0.05-0.5 \\ \hline
HUGO-MAT-02\tablefootnote{\href{https://spotmaterial.com/material/titanium/titanium-alloy-ti-6al-4v-grade-5-ti-6-4/}{Spot material} (accessed 11.05.2026)}              & 21-25 & 8.1-26.5 \\ \hline
HUGO-MAT-03\tablefootnote{\href{https://spotmaterial.com/material/aluminium/7000-series-alznmgcu-7075-7475/}{Spot material} (accessed 11.05.2026)}          & 2.55-5.1 & 4-20\\ \hline
HUGO-MAT-04\tablefootnote{\href{https://spotmaterial.com/material/magnesium/cast-we-alloys-we43-we54/}{Spot material} (accessed 11.05.2026)}               & 2.975-5.95 & 8.5-36 \\ \hline
HUGO-MAT-05\tablefootnote{\href{https://spotmaterial.com/material/alloy-steel/high-speed-tool-steel-hss-m-t/}{Spot material} (accessed 11.05.2026)}    & 8-20 & 6.8 \\ \hline
HUGO-MAT-06\tablefootnote{\href{https://spotmaterial.com/material/carbon-fibre-cf/high-modulus-hm-ep-cf60-60-continuous-cf-hm/}{Spot material} (accessed 11.05.2026)}    & 25,91-32.30  & 17.5-23.2 \\ \hline
HUGO-MAT-07\tablefootnote{\href{https://spotmaterial.com/material/high-modulus-glass-fibre/high-strength-s/}{Spot material} (accessed 12.05.2026)}            & 6.5-8.7  & 4.8 \\
\hline
\end{tabular}
\end{table}
\end{comment}


\subsection{Sustainability}
For the sustainability several factors had to be taken into account. There are the emissions stemming from the production and manufacture of the material and then there is the environmental impact of the material. The carbon footprint of a material is a method to quantify the amount of greenhouse emissions (measured in equivalent $CO_2$) necessary for the production of 1 kg of the material.   
To asses the environmental impact of a material at this stage it is mostly relevant to look at the different end-of-life scenarios for the material. To systematically do this it is necessary to perform a preliminary LCA-analyses of the aircraft using each chosen material. In this stage it makes sense to make a streamlined LCA where the focus will be on differences between the impact of the different materials\cite{Vogtlander2012}. This was performed for each material using the following model:   
\begin{figure}[H]
    \centering
    \includegraphics[width=1\linewidth]{figures/simpLCA-analysis.png}
    \caption{Simplified LCA model for the aircraft}
    \label{fig:placeholder}
\end{figure}
For the LCA-analysis the software openLCA was used with the added database: Swiss Federal Administration LCI Database (BAFU:2025) and the openLCA LCIA Methods package (insert reference here). For the streamlined assessment the focus was on the construction of fuselage, empennage and wings. The complete graphs and data files for each material can be found at the github (insert reference here)  





\subsection{Trade-off table}
\label{sub:materials_trade-off}
As previously mentioned the aircraft could be parted into several subsections. These subsections would each have their own requirements and trade-off tables. 


\begin{longtable}{|c|p{13cm}|}
\caption{Trade-Off Criteria for HUGO Configurations}
\label{tab:trade-off_criteria} \\

\hline
\textbf{Criterion} & \centering{\textbf{Rationale}} \\ 
\noalign
\hline\hline
\endfirsthead

\hline
\textbf{Criterion} & \textbf{Rationale} \\ 
\hline\hline
\endhead

\hline
\multicolumn{2}{r}{\textit{Continued on next page.}} \\
\endfoot

\hline
\endlastfoot

Structural Properties & The specific modulus ($E/\rho$), specific strength ($\sigma_y/\rho$), specific fracture toughness ($K_{Ic}/\rho$), column buckling efficiency ($\sqrt{E}/\rho$), and panel/skin buckling efficiency ($\sqrt[3]{E}/\rho$) are the most relevant properties to be analysed \cite{Arnold2012MaterialsSelection}. These shall be considered for the primary loading modes of each subsystem, as well as weight reduction.
\\ \hline

Durability & Based on HUGO's envisioned minimal operational life of 15 years (HUG-STR-012\footref{url:baseline_report}), the selected materials shall be capable of withstanding the design loads and environmental conditions over long periods of time, as well as multiple loading cycles within a day.
\\ \hline

Material Cost & Although raw material expenses constitute a fraction of the total unit cost compared to the processing, tooling, and labour required for low-volume manufacturing, they must still be evaluated to ensure the baseline economic viability of the selected candidates.
\\ \hline

Manufacturing Cost & Due to the low volume production model envisioned for HUGO (<10 units per batch), manufacturing costs are expected to drive the overall production cost of a single unit, and shall therefore be minimised.
\\ \hline

Design Freedom & Since HUGO has a high performance requirement, and its research and development costs are not constraining to the financial feasibility of the project, there is high value in choosing a material which allows for the production of complex shapes and profiles. A high design freedom is expected to increase the potential of HUGO as a flying testbed.
\\ \hline

Sustainability & While not being a strict design objective of HUGO, sustainability of the manufacturing of HUGO must be considered when choosing suitable materials. Standing at the forefront of aviation research, HUGO shall embrace the mission to lower the industry's environmental impact.
\\ \hline

\end{longtable}



\begin{longtable}{|p{0.24\textwidth}|| p{0.68\textwidth}|}
\caption{Trade-off Criteria for Material Selection}
\label{maintenance_operation} \\

\hline
\textbf{Maintenance Activity} & \textbf{Description} \\
\hline\hline
\endfirsthead

\hline
\textbf{Maintenance Activity} & \textbf{Description} \\
\hline\hline
\endhead

\hline
\multicolumn{2}{r}{\textit{Continued on next page.}} \\
\endfoot

\hline
\endlastfoot

Structural properties &  When evaluating different materials, the specific modulus ($E/\rho$), specific strength ($\sigma_y/\rho$), specific fracture toughness ($K_{Ic}/\rho$), column buckling efficiency ($\sqrt{E}/\rho$), and panel/skin buckling efficiency ($\sqrt[3]{E}/\rho$) are most relevant since weight reduction is crucial in aircraft design \cite{Arnold2012MaterialsSelection}. These properties are listed in \autoref{tab:material-properties} In order to systematically rank these choices based on the properties, the following method was used. First, each property was normalised, with 1.0 being the maximum value, and the rest scaled accordingly. Then, the properties were assigned a weight based on their relative importance. More details and explaination about this analysis are presented in \autoref{tab:Material_Tradeoff_Structural_Properties}.From this analysis, HUGO-MAT-06 (Carbon / Epoxy Composite) came out as the optimal material. However, glass fiber composite ranks very close as well, followed by sitka spruce, followed by the metal options.    
\\
\hline

Durability  & The durability of a specific material will be highly dependent on the use-case and operating-conditions of the aircraft. For HUGO-MAT-01 (spruce) humidity is a known major issue. For untreated wood, exposure to moisture leads to decomposing which greatly diminishes the structural properties of the material and can lead to warping. A coating would have to be applied to protect the wood. For metals, oxidation and corrosion is the largest issue. HUGO-MAT-02 (Ti) posses excellent corrosion and oxidation resistance. HUGO-MAT-02 also posses a good creep resistance \cite{Mridha2016}. For HUGO-MAT-03 (Al) its reaction with oxygen leads to a protective layer ($Al_2O_3$) which prevents further oxidation. The HUGO-MAT-03 would perform well in standard atmospheric conditions (primary area of operation for the aircraft). However, galvanic corrosion is an issue when it is contact with a more noble material \cite{Zhu2023}. Similarly, HUGO-MAT-04 (Mg) also forms a thin protective layer when in contact with oxygen, however it is the most susceptible to galvanic corrosion when in contact to more noble metals. The material also posses good creep-resistance \cite{Ouyang2020}. HUGO-MAT-05 (steel) is highly susceptible to corrosion and rusting. However, it posses a good resistance to creep and fatigue  \todo{insert citation here}. Both HUGO-MAT-06 (Carbon/epoxy comp.) and HUGO-MAT-07 (Glass/epoxy comp.) are also susceptible to oxidation, primarily caused due to UV exposure degrading the matrix. This can be prevented or reduced using protective layers or an opaque paint \cite{Composites_Corosion}. Water absorption also poses a problem in the long term \cite{Dong2024}. Fatigue also poses an issue for the composites. As damages from cyclically loads gradually builds up during operation the part might experience a apparently sudden catastrophic failure \cite{Mandell1981}.
 \\
\hline

Material cost & The cost of the raw materials of the structure of HUGO has a negligible overall impact on the financial feasibility of the project, as found in the market analysis. The cost of processing these materials and production of the parts have a far greater impact than the cost of the materials themselves. The specific costs of all the considered materials are listed in \autoref{tab:base-material-properties}.
\\
\hline

Manufacturing cost & The manufacturing cost is directly related to manufacturing complexity. Different designs can require different manufacturing processes. These processes can have different Non-Recurring Engineering costs (NRE). Due to the low volume nature of our project (<10 per production batch), the aim of this project should be to minimize these costs. A design comprising primarily of wood, as is the case with HUGO-MAT-01, the construction would likely comprise of machined spars, ribs and stiffeners, with a steam formed or bent laminate skin, bonded to the spars using adhesives. These techniques could allow for complex curves, with relatively low NRE costs for tooling and moulds. In the case of any of the metals chosen, (HUGO-MAT-02, HUGO-MAT-03, HUGO-MAT-04, HUGO-MAT-05), the construction would likely comprise of a thin, sheet metal skin, with stiffeners, ribs and spars fastened to the skin using fasteners, welded or bonded with an adhesive. If the skin has a simple bend, it can easily be laser cut and rolled to the correct radius. However, forming the sheet metal into complex, aerodynamic shapes would require a hydraulic sheet metal press with custom forms, or a stretch forming process, also requiring custom stretch blocks, both of which have high NRE costs. Incremental sheet forming and digital sheet forming techniques are beginning to gain traction, which would allow for complex shapes with no NRE costs, however this still has a low TRL and is still a niche technology. The ribs, spars and stiffeners would also have to be custom made, which would be expensive at the production volume of HUGO. Additionally, the wings, empennage and canards may have to be produced specifically for an experiment. For these cases, lead times and the cost of custom tooling for one-off parts would be prohibitively expensive. Composites require relatively inexpensive moulds which could be made of MDF or similar materials, and custom cut prepreg sheets for optimal fiber layup. This has a relatively low NRE cost. However, the layup process is laborious, time consuming and requires relatively high skill. After the layup, which could possibly involve an internal honecomb layer between the outer layers of fibre, the laminate needs to be bagged and cured in an autoclave for optimal curing, while avoiding voids and dry areas in the laminate. Autoclaves are often high demand equipment, which adds cost to the manufacturing process. However the economics change for parts such as the engine pylons, landing gear struts and other parts. These parts can simply be machined out of a block of metal. Although some weight can be saved by making these parts out of a composite material, each part would require a unique mould, with a different lay up technique and more autoclave time. 
\todo{citation needed}
 \\
\hline

Design Freedom  & Since HUGO has a high performance requirement, and since research and development costs are not constraining to the financial feasibility of the project, there is high value in choosing a material which allows for maximum design freedom. For instance, restricting the skin to only simple curves would restrict the aerodynamic design of the fuselage. Although HUGO-MAT-01 (Spruce) has similar strength to weight ratios to the metal options (\autoref{tab:material-properties}, it would require a very high volume of material to achieve the same strength. This would severely restrict our design freedom.
\\
\hline

Sustainability  &  The evaluation of the overall sustainability of a material is the combined consideration of several factors in the span of the materials lifetime. HUGO-MAT-01 has an extreme low carbon footprint which mostly stems from the extraction and  of since wood by it self extracts and binds carbon-dioxide from the atmosphere. Recycling HUGO-MAT-01 is also possible though mostly due to wear and tear down-cycling is most often the most viable approach. The material can be made into small chips which can then be compressed and made into panels\footnote{\href{https://www.unilinpanels.com/en/sustainability/recycling   }{UNILIN} (accessed 18.05.2026)}. Disposing HUGO-MAT-01 is also relatively simple as the material can be burned and used for energy and heating, though it is worth noting that this will be dependent on the surface coatings and treatment used on the material. Combustion will also lead to a release of the bounded carbon-dioxide. The majority of HUGO-MAT-02 (Ti) can be found in relative large quantities in the earth crust but the refinement and purification is a very energy-intensive process \cite{min11121425}. The recycle-ability of HUGO-MAT-02 is quite good and due to high production and refinement cost there is a clear incentive for recycling \cite{RAHIMI2025107056}. The production of both HUGO-MAT-03 and HUGO-MAT-04 has a large carbon footprint (see \autoref{tab:material-economics} with HUGO-MAT-04 being very high due to the specific refining process used\cite{dot}. HUGO-MAT-05 has a medium carbon footprint. As is often the case with metals it is possible for the recycled HUGO-MAT-02,  HUGO-MAT-03,  HUGO-MAT-04 and HUGO-MAT-05 to retain most of the structural performance of the virgin material \cite{euric_metal_recycling}. The final two materials would be HUGO-MAT-06 and HUGO-MAT-07. Sustainability is often presented as one of the largest issues for the broad implementation for carbon and glass-composites \cite{dot}. The production of both fibres are  extremely energy-heavy see \autoref{tab:material-economics} with especially carbon fibre having a large carbon footprint. There lies a potential sourcing the carbon from the atmosphere though this is not implemented on a industrial scale \cite{Keiner2024}. The resin used for the matrix is typically a thermoset. Recycling of these materials is typically very difficult and the large majority of glass and carbon- composites ends in landfills or are incinerated \cite{dot}.         
 \\
\hline

\end{longtable}


\begin{table}[H]
\centering
\caption{Trade-Off Matrix for fuselage, wings and empennage}
\renewcommand{\arraystretch}{1.4}

\begin{adjustbox}{width=\textwidth}
\begin{tabular}{|
p{3cm}||
>{\centering\arraybackslash}p{3cm}|
>{\centering\arraybackslash}p{3cm}|
>{\centering\arraybackslash}p{3cm}|
>{\centering\arraybackslash}p{3cm}|
>{\centering\arraybackslash}p{3cm}|
>{\centering\arraybackslash}p{3cm}||
>{\centering\arraybackslash}p{3cm}|
}
\hline
\rowcolor[HTML]{4472C4}
\color{white}\textbf{Material} &
\color{white}\textbf{Structural properties} &
\color{white}\textbf{Durability} &
\color{white}\textbf{Material cost} &
\color{white}\textbf{Manufacturing cost} &
\color{white}\textbf{Design freedom} &
\color{white}\textbf{\shortstack{Sustain-\\ability}} &
\color{white}\textbf{Weighted Total} \\
\hline\hline

\rowcolor[HTML]{D9E2F3}
\textbf{Weight} &
1 &
0.85 &
0.1 &
0.5 &
1 &
0.5 &
\\
\hline

HUGO-MAT-01
& 5 \cellcolor[HTML]{63BE7B}
& 1 \cellcolor[HTML]{F8696B}
& 5 \cellcolor[HTML]{63BE7B}
& 3 \cellcolor[HTML]{F4E285}
& 1 \cellcolor[HTML]{F8696B}
& 5 \cellcolor[HTML]{63BE7B}
& 11.35 \\
\hline

\rowcolor[HTML]{D9E2F3}
HUGO-MAT-02
& 3 \cellcolor[HTML]{F4E285}
& 5 \cellcolor[HTML]{63BE7B}
& 1 \cellcolor[HTML]{F8696B}
& 1 \cellcolor[HTML]{F8696B}
& 3 \cellcolor[HTML]{F4E285}
& 2 \cellcolor[HTML]{F4B183}
& 11.85 \\
\hline

HUGO-MAT-03
& 3 \cellcolor[HTML]{F4E285}
& 4 \cellcolor[HTML]{A9D18E}
& 4 \cellcolor[HTML]{A9D18E}
& 2 \cellcolor[HTML]{F4B183}
& 3 \cellcolor[HTML]{F4E285}
& 3 \cellcolor[HTML]{F4E285}
& 12.30 \\
\hline

\rowcolor[HTML]{D9E2F3}
HUGO-MAT-04
& 3 \cellcolor[HTML]{F4E285}
& 3 \cellcolor[HTML]{F4E285}
& 4 \cellcolor[HTML]{A9D18E}
& 1 \cellcolor[HTML]{F8696B}
& 3 \cellcolor[HTML]{F4E285}
& 2 \cellcolor[HTML]{F4B183}
& 10.45 \\
\hline

HUGO-MAT-05
& 3 \cellcolor[HTML]{F4E285}
& 2 \cellcolor[HTML]{F4B183}
& 1 \cellcolor[HTML]{F8696B}
& 2 \cellcolor[HTML]{F4B183}
& 3 \cellcolor[HTML]{F4E285}
& 3 \cellcolor[HTML]{F4E285}
& 10.30 \\
\hline

\rowcolor[HTML]{D9E2F3}
\textbf{HUGO-MAT-06}
& \textbf{5} \cellcolor[HTML]{63BE7B}
& \textbf{4} \cellcolor[HTML]{A9D18E}
& \textbf{2} \cellcolor[HTML]{F4B183}
& \textbf{2} \cellcolor[HTML]{F4B183}
& \textbf{5} \cellcolor[HTML]{63BE7B}
& \textbf{1} \cellcolor[HTML]{F8696B}
& \textbf{15.10} \\
\hline

HUGO-MAT-07
& 5 \cellcolor[HTML]{63BE7B}
& 4 \cellcolor[HTML]{A9D18E}
& 4 \cellcolor[HTML]{A9D18E}
& 2 \cellcolor[HTML]{F4B183}
& 4 \cellcolor[HTML]{A9D18E}
& 1 \cellcolor[HTML]{F8696B}
& 14.30 \\
\hline

\end{tabular}
\end{adjustbox}
\label{tab:fuselage_trade-off_matrix_material_selection}
\end{table}


\begin{table}[H]
\centering
\caption{Trade-Off Matrix for Landing gear and Pylons}
\renewcommand{\arraystretch}{1.4}

\begin{adjustbox}{width=\textwidth}
\begin{tabular}{|
p{3cm}||
>{\centering\arraybackslash}p{3cm}|
>{\centering\arraybackslash}p{3cm}|
>{\centering\arraybackslash}p{3cm}|
>{\centering\arraybackslash}p{3cm}|
>{\centering\arraybackslash}p{3cm}|
>{\centering\arraybackslash}p{3cm}||
>{\centering\arraybackslash}p{3cm}|
}
\hline
\rowcolor[HTML]{4472C4}
\color{white}\textbf{Material} &
\color{white}\textbf{Structural properties} &
\color{white}\textbf{Durability} &
\color{white}\textbf{Material cost} &
\color{white}\textbf{Manufacturing cost} &
\color{white}\textbf{Design freedom} &
\color{white}\textbf{\shortstack{Sustain-\\ability}} &
\color{white}\textbf{Weighted Total} \\
\hline\hline

\rowcolor[HTML]{D9E2F3}
\textbf{Weight} &
1 &
0.85 &
0.1 &
0.5 &
1 &
0.5 &
\\
\hline

HUGO-MAT-01
& 5 \cellcolor[HTML]{63BE7B}
& 1 \cellcolor[HTML]{F8696B}
& 5 \cellcolor[HTML]{63BE7B}
& 5 \cellcolor[HTML]{63BE7B}
& 1 \cellcolor[HTML]{F8696B}
& 5 \cellcolor[HTML]{63BE7B}
& 12.35 \\
\hline

\rowcolor[HTML]{D9E2F3}
HUGO-MAT-02
& 3 \cellcolor[HTML]{F4E285}
& 5 \cellcolor[HTML]{63BE7B}
& 1 \cellcolor[HTML]{F8696B}
& 3 \cellcolor[HTML]{F4E285}
& 5 \cellcolor[HTML]{63BE7B}
& 2 \cellcolor[HTML]{F4B183}
& 14.85 \\
\hline

\textbf{HUGO-MAT-03}
& \textbf{3} \cellcolor[HTML]{F4E285}
& \textbf{4} \cellcolor[HTML]{A9D18E}
& \textbf{4} \cellcolor[HTML]{A9D18E}
& \textbf{4} \cellcolor[HTML]{A9D18E}
& \textbf{5} \cellcolor[HTML]{63BE7B}
& \textbf{3} \cellcolor[HTML]{F4E285}
& \textbf{15.30} \\
\hline

\rowcolor[HTML]{D9E2F3}
HUGO-MAT-04
& 3 \cellcolor[HTML]{F4E285}
& 3 \cellcolor[HTML]{F4E285}
& 4 \cellcolor[HTML]{A9D18E}
& 3 \cellcolor[HTML]{F4E285}
& 5 \cellcolor[HTML]{63BE7B}
& 2 \cellcolor[HTML]{F4B183}
& 13.45 \\
\hline

HUGO-MAT-05
& 3 \cellcolor[HTML]{F4E285}
& 2 \cellcolor[HTML]{F4B183}
& 1 \cellcolor[HTML]{F8696B}
& 4 \cellcolor[HTML]{A9D18E}
& 5 \cellcolor[HTML]{63BE7B}
& 3 \cellcolor[HTML]{F4E285}
& 13.30 \\
\hline

\rowcolor[HTML]{D9E2F3}
HUGO-MAT-06
& 5 \cellcolor[HTML]{63BE7B}
& 4 \cellcolor[HTML]{A9D18E}
& 2 \cellcolor[HTML]{F4B183}
& 1 \cellcolor[HTML]{F8696B}
& 3 \cellcolor[HTML]{F4E285}
& 1 \cellcolor[HTML]{F8696B}
& 12.60 \\
\hline

HUGO-MAT-07
& 5 \cellcolor[HTML]{63BE7B}
& 4 \cellcolor[HTML]{A9D18E}
& 4 \cellcolor[HTML]{A9D18E}
& 1 \cellcolor[HTML]{F8696B}
& 2 \cellcolor[HTML]{F4B183}
& 1 \cellcolor[HTML]{F8696B}
& 11.80 \\
\hline

\end{tabular}
\end{adjustbox}
\label{tab:landing_gear_trade-off_matrix_material_selection}
\end{table}



\section{Sustainable Development Strategy}
To access the environmental impact of the system, a preliminary LCA-analysis was set up using the open-source software openLCA. At this point it 

\section{Sensitivity study}